\documentclass[11pt,oneside]{book}

\usepackage[T1]{fontenc}
\usepackage[utf8]{inputenc}
\usepackage{lmodern}
\usepackage[a4paper,margin=25mm]{geometry}
\usepackage{microtype}
\usepackage{amsmath,amssymb,mathtools}
\usepackage{booktabs}
\usepackage{longtable}
\usepackage{array}
\usepackage{enumitem}
\usepackage{xcolor}
\usepackage{listings}
\usepackage{fancyhdr}
\usepackage{hyperref}
\usepackage[nameinlink,noabbrev]{cleveref}

\hypersetup{
  colorlinks=true,
  linkcolor=blue!55!black,
  urlcolor=blue!65!black,
  citecolor=blue!55!black,
  pdftitle={DiscoGMPI User Guide},
  pdfauthor={DiscoGMPI project}
}

\pagestyle{fancy}
\fancyhf{}
\fancyhead[L]{DiscoGMPI User Guide}
\fancyhead[R]{\leftmark}
\fancyfoot[C]{\thepage}
\setlength{\headheight}{14pt}

\definecolor{codebg}{RGB}{247,247,247}
\definecolor{codeframe}{RGB}{210,210,210}
\definecolor{codegreen}{RGB}{25,115,45}
\definecolor{codeblue}{RGB}{15,70,145}
\definecolor{codepurple}{RGB}{120,45,135}

\lstdefinelanguage{Julia}{
  morekeywords={
    abstract,break,catch,const,continue,do,else,elseif,end,export,false,
    finally,for,function,global,if,import,in,let,local,macro,module,mutable,
    primitive,quote,return,struct,true,try,using,where,while
  },
  sensitive=true,
  morecomment=[l]\#,
  morestring=[b]",
  morestring=[b]'
}

\lstset{
  language=Julia,
  basicstyle=\ttfamily\small,
  keywordstyle=\color{codeblue}\bfseries,
  commentstyle=\color{codegreen},
  stringstyle=\color{codepurple},
  backgroundcolor=\color{codebg},
  frame=single,
  rulecolor=\color{codeframe},
  breaklines=true,
  breakatwhitespace=false,
  showstringspaces=false,
  columns=fullflexible,
  keepspaces=true,
  tabsize=4,
  numbers=left,
  numberstyle=\tiny\color{gray},
  numbersep=8pt,
  xleftmargin=2em,
  framexleftmargin=1.5em,
  aboveskip=0.8em,
  belowskip=0.8em,
  literate=
    {ε}{{$\varepsilon$}}1
    {μ}{{$\mu$}}1
    {Δ}{{$\Delta$}}1
    {π}{{$\pi$}}1
}

\newcommand{\pkg}{\texttt{DiscoGMPI}}
\newcommand{\code}[1]{\texttt{#1}}
\newcommand{\api}[1]{\path{#1}}
\newcommand{\vect}[1]{\boldsymbol{#1}}
\newcommand{\jump}[1]{\left[\!\left[#1\right]\!\right]}
\newcommand{\dd}{\,\mathrm{d}}
\setlength{\emergencystretch}{2em}

\title{\Huge DiscoGMPI\\[0.5em]\Large User Guide and Technical Reference}
\author{DiscoGMPI project}
\date{Version 0.1.0\\June 2026}

\begin{document}

\frontmatter
\hypersetup{pageanchor=false}
\maketitle
\hypersetup{pageanchor=true}
\setcounter{page}{1}

\chapter*{About this guide}
\addcontentsline{toc}{chapter}{About this guide}

This manual describes the \pkg{} software as implemented in the repository at
the time of writing. It covers the mathematical model, discontinuous Galerkin
(DG) discretization, Maxwell formulations, MPI decomposition, installation,
mesh preparation, public workflows, examples, tests, benchmarks, diagnostics,
and known limitations.

\pkg{} is the distributed-memory development line derived from DiscoG3D. It is
a separate Julia package and repository. Work performed in \pkg{} does not
overwrite the original DiscoG3D source tree.

\paragraph{Scope.}
The principal production path documented here is the source-free,
three-dimensional Maxwell system on affine tetrahedral meshes, with
perfect-electric-conductor (PEC) boundaries and explicit time integration.
The serial and threaded solver also contains periodic-boundary infrastructure.
The MPI Maxwell driver currently exchanges inter-rank face traces and supports
PEC boundaries, but distributed periodic boundaries, sources, perfectly
matched layers (PML), and heterogeneous material jumps remain future work.

\paragraph{Notation.}
Superscripts $-$ and $+$ denote the local and neighboring traces at a face.
The unit normal $\vect n$ points out of the minus element. The polynomial
degree is denoted by $N$, the number of volume nodes per tetrahedron by
$N_p$, and the number of nodes per triangular face by $N_{fp}$.

\tableofcontents
\listoftables

\mainmatter

\chapter{Overview}

\section{Purpose}

\pkg{} provides:

\begin{itemize}[leftmargin=*]
  \item nodal DG operators for affine tetrahedral meshes;
  \item source-free Maxwell volume and surface operators;
  \item Hesthaven--Warburton strong-form and Poisson-bracket formulations;
  \item central and upwind Maxwell fluxes where mathematically supported;
  \item serial and Julia-threaded local execution backends;
  \item explicit Runge--Kutta methods of orders one through five;
  \item explicit partitioned symplectic Runge--Kutta (ESPRK) schemes of
        orders one through six;
  \item METIS element partition input;
  \item owned-element and one-layer face-halo mesh construction;
  \item MPI face orientation and nodal permutation handling;
  \item nonblocking exchange of the six Maxwell traces;
  \item distributed energy reductions and distributed time integration;
  \item partition visualization, regression tests, convergence cases, and
        backend benchmarks.
\end{itemize}

\section{Repository layout}

\begin{longtable}{p{0.39\textwidth}p{0.53\textwidth}}
\toprule
Path & Purpose \\
\midrule
\endhead
\api{src/DiscoGMPI.jl} & Package module, includes, and exported API. \\
\api{src/solver/} & Serial and threaded DG/Maxwell data structures and kernels. \\
\api{src/DistributedMesh3D.jl} & Ownership, ghost elements, local/global maps,
face adjacency, and neighbor communication tables. \\
\api{src/DistributedDG.jl} & Rank-local DG construction, root-only mesh
distribution, Maxwell trace exchange, distributed right-hand side, energy,
and time integration. \\
\api{src/TraceMaps.jl} & General MPI trace maps and face-node permutation support. \\
\api{src/MetisIO.jl} & Legacy VTK tetrahedral reader and METIS mesh writer. \\
\api{src/PartitionVTK.jl} & Legacy VTK partition visualization writers. \\
\api{examples/solver/} & Serial and threaded Maxwell examples and convergence
drivers inherited from DiscoG3D. \\
\api{examples/distributed_maxwell.jl} & End-to-end MPI Maxwell demonstration. \\
\api{examples/meshes/} & Small sample mesh and METIS partition files. \\
\api{benchmark/solver/} & Serial-versus-threaded right-hand-side benchmarks and
plotting scripts. \\
\api{test/} & Package tests and explicitly launched MPI tests. \\
\bottomrule
\end{longtable}

\section{Supported execution modes}

\begin{description}[style=nextline]
  \item[Serial]
  One Julia process using \code{SerialBackend()}.

  \item[Threaded]
  One Julia process using \code{ThreadedBackend()} and a Julia thread count
  selected with \code{--threads=N} or \code{JULIA\_NUM\_THREADS=N}.

  \item[Distributed memory]
  Multiple Julia processes launched with \code{mpiexec}. Every rank owns a
  subset of tetrahedra and stores a one-face halo.

  \item[Hybrid MPI plus threads]
  The distributed constructor accepts a backend, so a threaded local backend
  can be selected per rank. This path should be validated on the target
  machine before production use; the primary MPI regression tests use the
  serial local backend.
\end{description}

\chapter{Mathematical Model}

\section{Source-free Maxwell equations}

For scalar, positive, spatially uniform permittivity $\varepsilon$ and
permeability $\mu$, \pkg{} advances
\begin{align}
  \varepsilon \frac{\partial \vect E}{\partial t}
    &= \nabla \times \vect H, \label{eq:maxwell-e}\\
  \mu \frac{\partial \vect H}{\partial t}
    &= -\nabla \times \vect E. \label{eq:maxwell-h}
\end{align}
The wave speed, impedance, and admittance are
\begin{equation}
  c = \frac{1}{\sqrt{\varepsilon\mu}},
  \qquad
  Z = \sqrt{\frac{\mu}{\varepsilon}},
  \qquad
  Y = \frac{1}{Z}.
\end{equation}

The divergence constraints
\begin{equation}
  \nabla\cdot(\varepsilon\vect E)=0,
  \qquad
  \nabla\cdot(\mu\vect H)=0
\end{equation}
are not solved as separate equations. Their preservation depends on the
initial condition, mesh, spatial discretization, flux, and time integrator.

\section{Electromagnetic energy}

The continuous energy is
\begin{equation}
  \mathcal E(t)
  =
  \frac{1}{2}\int_\Omega
  \left(
    \varepsilon |\vect E|^2 + \mu |\vect H|^2
  \right)\dd \vect x.
\end{equation}
The code evaluates its DG approximation element by element using the reference
mass matrix and affine Jacobian determinant:
\begin{equation}
  \mathcal E_h
  =
  \frac{1}{2}
  \sum_{K}
  |J_K|
  \left[
    \varepsilon\sum_{\alpha\in\{x,y,z\}}
      \vect E_{\alpha,K}^{T}M\vect E_{\alpha,K}
    +
    \mu\sum_{\alpha\in\{x,y,z\}}
      \vect H_{\alpha,K}^{T}M\vect H_{\alpha,K}
  \right].
\end{equation}
In MPI runs only owned elements contribute locally, followed by an
\code{MPI.Allreduce} sum. Ghost elements are therefore not double counted.

\section{PEC boundary condition}

The perfect-electric-conductor condition is
\begin{equation}
  \vect n\times\vect E = \vect 0
  \qquad\text{on }\partial\Omega_{\mathrm{PEC}}.
\end{equation}
\pkg{} imposes it through a reflected exterior trace:
\begin{align}
  \vect E^+ &= -\vect E^- + 2(\vect n\cdot\vect E^-)\vect n,\\
  \vect H^+ &= \vect H^-.
\end{align}
This reverses the tangential electric field while retaining its normal
component. Boundary-condition selection is tag based. A typical registry is:

\begin{lstlisting}
registry = MaxwellBoundaryRegistry(
    Dict(10 => DiscoGMPI.MaxwellBC_PEC),
)
\end{lstlisting}

An unregistered boundary tag maps to \code{MaxwellBC\_None}; it does not
silently become PEC.

\chapter{Tetrahedral DG Discretization}

\section{Mesh and approximation space}

The physical domain is partitioned into affine tetrahedra $K$. On each element,
every scalar field component is represented in
\begin{equation}
  \mathbb P^N(K)
  =
  \{p(x,y,z): \deg p\leq N\},
\end{equation}
with
\begin{equation}
  N_p = \frac{(N+1)(N+2)(N+3)}{6}
\end{equation}
volume nodes. A triangular face contains
\begin{equation}
  N_{fp} = \frac{(N+1)(N+2)}{2}
\end{equation}
nodes.

The reference tetrahedron has vertices
\[
(-1,-1,-1),\quad
(1,-1,-1),\quad
(-1,1,-1),\quad
(-1,-1,1).
\]
The implementation builds equispaced nodal points, an orthonormal polynomial
basis, a Vandermonde matrix, nodal derivative matrices, an exact polynomial
mass matrix, and face mass/lifting operators.

\section{Affine mapping and physical derivatives}

For an affine map $\vect x=\vect x_K(r,s,t)$, let $J_K$ denote its Jacobian.
Reference gradients are transformed by
\begin{equation}
  \nabla_{\vect x} u
  =
  J_K^{-T}\nabla_{(r,s,t)}u.
\end{equation}
The physical derivative matrices $D_x,D_y,D_z$ and weak derivative matrices
\begin{equation}
  S_x = M D_x,\qquad S_y=M D_y,\qquad S_z=M D_z
\end{equation}
are precomputed for each tetrahedron. The Poisson-bracket formulation also
uses $S_x^T,S_y^T,S_z^T$.

\section{Topology, traces, and lifting}

Faces are identified by sorted triples of global vertex IDs. A face appearing
once is a boundary face; a face appearing twice is an interior face. More than
two adjacent tetrahedra is treated as a non-manifold error.

For every interior face, the DG trace map stores:
\begin{itemize}[leftmargin=*]
  \item minus and plus element IDs;
  \item local face IDs;
  \item face-node indices;
  \item the plus-to-minus nodal permutation;
  \item the outward minus-side normal and face area.
\end{itemize}

Surface residuals are projected through the face mass matrix and lifted into
the element volume. Conceptually, a face correction $\vect f$ contributes
\begin{equation}
  M_K^{-1}M_{f,K}\vect f
\end{equation}
to the nodal right-hand side, with the physical face area and element Jacobian
included by the implementation.

\section{Hesthaven--Warburton formulation}

The strong-form volume terms directly approximate
\begin{align}
  \dot{\vect E}_K &= \frac{1}{\varepsilon}
    (\nabla_h\times\vect H)_K + \text{surface correction},\\
  \dot{\vect H}_K &= -\frac{1}{\mu}
    (\nabla_h\times\vect E)_K + \text{surface correction}.
\end{align}

Let $\jump{\vect E}=\vect E^+-\vect E^-$ and
$\jump{\vect H}=\vect H^+-\vect H^-$. The central correction used by the
code is
\begin{align}
  \vect f_E^{\mathrm c}
    &= \frac{1}{2}\vect n\times\jump{\vect H},\\
  \vect f_H^{\mathrm c}
    &= -\frac{1}{2}\vect n\times\jump{\vect E}.
\end{align}
The upwind flux adds tangential penalties:
\begin{align}
  \vect f_E^{\mathrm up}
    &= \vect f_E^{\mathrm c}
       -\frac{Y}{2}\vect n\times
        \left(\vect n\times\jump{\vect E}\right),\\
  \vect f_H^{\mathrm up}
    &= \vect f_H^{\mathrm c}
       -\frac{Z}{2}\vect n\times
        \left(\vect n\times\jump{\vect H}\right).
\end{align}

Select the formulation with:

\begin{lstlisting}
central = HesthavenWarburtonFormulation(MaxwellFlux_Central)
upwind  = HesthavenWarburtonFormulation(MaxwellFlux_Upwind)
\end{lstlisting}

The central flux is nondissipative in the ideal compatible setting. The
upwind flux introduces numerical dissipation and is normally preferred when
robust damping of unresolved jumps is more important than exact semidiscrete
energy conservation.

\section{Poisson-bracket formulation}

The Poisson-bracket implementation is a partitioned weak-operator
discretization. The electric update uses weak derivative matrices
$S_x,S_y,S_z$, while the magnetic update uses their transposes. This pairing
constructs a discrete skew structure and supports energy-compatible
partitioned time stepping.

For an interior face, the implemented corrections can be summarized as
\begin{align}
  \vect f_E^{\mathrm{PB}}
    &= -\frac{1}{2\varepsilon}
       \vect n\times(\vect H^- - \vect H^+),\\
  \vect f_H^{\mathrm{PB}}
    &= -\frac{1}{2\mu}
       \vect n\times(\vect E^- + \vect E^+).
\end{align}

Only the central flux is supported:

\begin{lstlisting}
formulation = PoissonBracketFormulation()
\end{lstlisting}

Attempting to combine the Poisson-bracket formulation with the upwind flux
raises an error because the upwind penalty is dissipative and does not define
the implemented partitioned Poisson operator.

\chapter{Time Integration}

\section{Explicit Runge--Kutta methods}

For a semidiscrete system $\dot U=R(U)$, an explicit Runge--Kutta method uses
\begin{align}
  U^{(i)} &= U^n + \Delta t\sum_{j<i}a_{ij}K_j,\\
  K_i &= R(U^{(i)}),\\
  U^{n+1} &= U^n + \Delta t\sum_i b_iK_i.
\end{align}

\begin{table}[htbp]
\centering
\caption{Explicit Runge--Kutta schemes available through
\code{explicit\_rk\_scheme(order)}.}
\begin{tabular}{clc}
\toprule
Order & Scheme & Stages \\
\midrule
1 & Forward Euler & 1 \\
2 & Explicit midpoint & 2 \\
3 & Classical third order & 3 \\
4 & Classical fourth order & 4 \\
5 & Dormand--Prince fifth-order update & 7 \\
\bottomrule
\end{tabular}
\end{table}

The Dormand--Prince implementation uses the fifth-order update only. Adaptive
step-size control and the embedded fourth-order estimator are not exposed.

\section{Explicit partitioned symplectic methods}

The Maxwell state is divided into electric and magnetic partitions. At every
ESPRK stage, one partition is updated using its weight, the right-hand side is
recomputed, and the complementary partition is updated using its own weight.
Orders one through six are available:

\begin{lstlisting}
scheme = explicit_partitioned_symplectic_rk_scheme(
    4;
    first_partition = :H,
)
\end{lstlisting}

For \code{PoissonBracketFormulation}, the magnetic partition must be advanced
first:

\begin{lstlisting}
formulation = PoissonBracketFormulation()

run_maxwell_partitioned_symplectic_time_steps!(
    U,
    dg,
    registry,
    formulation;
    psrk_order = 4,
    first_partition = :H,
    dt = dt,
    nsteps = nsteps,
)
\end{lstlisting}

Using the conventional explicit Runge--Kutta driver with the Poisson-bracket
formulation is rejected by the high-level API.

\section{CFL estimate}

The code defines the element characteristic length
\begin{equation}
  h_K = \frac{3V_K}{A_K},
\end{equation}
where $V_K$ is the tetrahedron volume and $A_K$ its total surface area. The
recommended estimate is
\begin{equation}
  \Delta t
  =
  \mathrm{CFL}
  \frac{\min_K h_K}{(2N+1)c}.
\end{equation}
It is obtained with:

\begin{lstlisting}
dt, sizes = estimate_maxwell_dt(
    mesh,
    dg.geometry,
    dg.ref;
    CFL = 0.10,
)
\end{lstlisting}

This expression is an estimate, not a proof of stability for every mesh,
formulation, flux, and RK scheme. Production studies should include a
step-size sensitivity check.

\chapter{Distributed-Memory Design}

\section{Element ownership and face halos}

An element partition vector assigns every global tetrahedron to a zero-based
MPI rank:
\begin{equation}
  p_e\in\{0,\ldots,P-1\}.
\end{equation}
For rank $r$, \pkg{} stores:

\begin{itemize}[leftmargin=*]
  \item every element for which $p_e=r$ as \emph{owned};
  \item every off-rank tetrahedron sharing a complete face with an owned
        element as a \emph{ghost};
  \item the nodes required by those owned and ghost elements;
  \item global-to-local and local-to-global maps;
  \item ownership, material IDs, face neighbors, volumes, and barycenters;
  \item one communication record per neighboring rank.
\end{itemize}

Only a one-face halo is needed because the Maxwell DG surface operator couples
elements only through face traces.

\section{Face orientation}

A triangular face admits six vertex permutations. For each inter-rank face,
the distributed mesh records the orientation code that maps the neighbor
ordering to the local ordering. The trace layer converts face nodes to
barycentric coordinates and precomputes one nodal permutation for each of the
six orientations. Received plus traces are therefore presented in the same
physical point ordering as the local minus trace.

This permutation step is essential for polynomial order $N>1$: matching only
the three vertices is insufficient because edge and interior face nodes must
also be reordered.

\section{Maxwell communication}

At each right-hand-side evaluation:

\begin{enumerate}[leftmargin=*]
  \item The six owned face traces
        $(E_x,E_y,E_z,H_x,H_y,H_z)$ are packed for every neighbor.
  \item All nonblocking receives are posted.
  \item All nonblocking sends are posted.
  \item \code{MPI.Waitall} completes the exchange.
  \item Received values are unpacked into the corresponding ghost face nodes.
  \item The local DG right-hand side is evaluated.
  \item Every ghost-element right-hand-side column is set to zero, ensuring
        that time integration updates owned elements only.
\end{enumerate}

For a neighbor with $F_{rn}$ interface faces, the message contains
\begin{equation}
  6N_{fp}F_{rn}
\end{equation}
double-precision values per right-hand-side evaluation.

\paragraph{Current computation model.}
The local DG kernel currently evaluates owned and stored ghost elements before
zeroing ghost residuals. Communication volume is already trace-only, but a
future optimization can avoid unnecessary ghost volume work and overlap
communication with owned volume calculations.

\section{Root-only mesh distribution}

The scalable constructor keeps the global mesh and partition vector only on a
designated root rank. Root builds each rank's owned-plus-halo package and sends
that serialized package to the destination. Non-root ranks never construct or
retain the complete global mesh.

\begin{lstlisting}
using MPI
using DiscoGMPI

MPI.Init()
comm = MPI.COMM_WORLD
rank = MPI.Comm_rank(comm)
nranks = MPI.Comm_size(comm)

root_mesh = rank == 0 ? read_vtu_mesh("mesh.vtu") : nothing
root_partition = if rank == 0
    read_metis_epart("mesh.mesh.epart.$nranks")
else
    nothing
end

ddg = build_distributed_dg_from_root(
    root_mesh,
    root_partition,
    2;
    comm = comm,
    root = 0,
)

MPI.Barrier(comm)
MPI.Finalize()
\end{lstlisting}

The file-based overload performs both reads only on root:

\begin{lstlisting}
ddg = build_distributed_dg_from_root(
    "mesh.vtu",
    "mesh.mesh.epart.4",
    2;
    comm = MPI.COMM_WORLD,
)
\end{lstlisting}

All ranks must call the constructor collectively. Root validation errors are
broadcast so that malformed input fails collectively rather than leaving
worker ranks blocked.

\section{Replicated constructor}

When a global \code{RawVTUMesh} and partition vector already exist on every
rank, the compatibility constructor remains available:

\begin{lstlisting}
ddg = build_distributed_dg(
    global_mesh,
    elem_to_rank,
    order;
    comm = MPI.COMM_WORLD,
)
\end{lstlisting}

The root-only constructor is preferable for memory scaling on non-root ranks.
Root still holds the complete input mesh during distribution.

\chapter{Installation and Environment}

\section{Prerequisites}

\begin{itemize}[leftmargin=*]
  \item Julia 1.9 or later.
  \item An MPI implementation compatible with the MPI library selected by
        MPI.jl.
  \item An \code{mpiexec} command for multi-process runs.
  \item METIS and its \code{mpmetis} executable when generating new
        partitions.
  \item ParaView or another VTK viewer for visualization.
  \item Optional: Python 3 and Matplotlib for benchmark plots.
  \item Optional: a LaTeX installation to build this manual.
\end{itemize}

\section{Instantiate the Julia environment}

Run all commands from the repository root:

\begin{lstlisting}[language=bash]
julia --project=. -e 'using Pkg; Pkg.instantiate()'
julia --project=. -e 'using DiscoGMPI; println("DiscoGMPI loaded")'
\end{lstlisting}

The package dependencies are recorded in \code{Project.toml} and
\code{Manifest.toml}. Important dependencies include MPI.jl, ReadVTK.jl,
VTKBase.jl, WriteVTK.jl, BenchmarkTools.jl, and Julia's LinearAlgebra standard
library.

\section{Check MPI}

\begin{lstlisting}[language=bash]
mpiexec -n 2 julia --project=. -e '
using MPI
MPI.Init()
println("rank ", MPI.Comm_rank(MPI.COMM_WORLD),
        " of ", MPI.Comm_size(MPI.COMM_WORLD))
MPI.Barrier(MPI.COMM_WORLD)
MPI.Finalize()
'
\end{lstlisting}

If this command fails before \pkg{} is loaded, fix the MPI.jl/runtime
configuration first. In particular, do not mix an arbitrary system
\code{mpiexec} with a different MPI library selected by MPI.jl.

\section{Thread configuration}

Thread counts are selected when Julia starts:

\begin{lstlisting}[language=bash]
julia --project=. --threads=8 my_case.jl
JULIA_NUM_THREADS=8 julia --project=. my_case.jl
\end{lstlisting}

Inside Julia, verify the setting with:

\begin{lstlisting}
println(Threads.nthreads())
\end{lstlisting}

\chapter{Mesh Preparation}

\section{RawVTUMesh format}

The primary solver mesh container is:

\begin{lstlisting}
RawVTUMesh(
    points,        # 3 x number_of_points, Float64
    tets,          # 4 x number_of_tetrahedra, Int
    tris,          # 3 x number_of_boundary_triangles, Int
    tet_cell_ids,
    tri_cell_ids,
    cell_data,
)
\end{lstlisting}

Connectivity uses Julia's one-based indexing. The XML VTU reader accepts
tetrahedral volume cells and triangular surface cells and ignores unsupported
cell types.

\section{Boundary and material tags}

By default, \code{DGDiscretization} looks for cell data named
\code{boundary\_id}. The array is indexed by the original VTK cell IDs;
\code{tri\_data(mesh, "boundary\_id")} extracts its surface-triangle values.

The distributed constructors also accept:

\begin{lstlisting}
boundary_tag_name = "boundary_id"
material_tag_name = "region_id"
\end{lstlisting}

Material IDs are preserved in the distributed topology. The current Maxwell
operator nevertheless uses scalar run-wide values of $\varepsilon$ and $\mu$;
element-dependent material jumps are not yet implemented.

\section{Inspect a VTU mesh}

\begin{lstlisting}[language=bash]
julia --project=. examples/solver/read_vtu_mesh.jl path/to/mesh.vtu
\end{lstlisting}

The inspection example prints mesh dimensions, connectivity, tags, topology,
geometry, mappings, basis diagnostics, face operators, trace checks, flux-face
checks, and several local operator tests. It is intentionally verbose and is a
useful first check for a new mesh.

\section{Generate a METIS mesh}

The sample workflow begins with a legacy ASCII VTK tetrahedral mesh:

\begin{lstlisting}[language=bash]
julia --project=. examples/write_metis.jl
\end{lstlisting}

This calls:

\begin{lstlisting}
write_metis_mesh_from_vtk(
    "examples/meshes/tet_mesh.vtk",
    "examples/meshes/tet_mesh.mesh",
)
\end{lstlisting}

The generated METIS mesh contains one tetrahedron per line. Create $P$
partitions with:

\begin{lstlisting}[language=bash]
mpmetis examples/meshes/tet_mesh.mesh 4
\end{lstlisting}

METIS normally writes:

\begin{itemize}[leftmargin=*]
  \item \code{tet\_mesh.mesh.epart.4}: one element partition ID per line;
  \item \code{tet\_mesh.mesh.npart.4}: one node partition ID per line.
\end{itemize}

\pkg{} reads the element partition file. Partition IDs are expected to be
zero based. For one-based files use:

\begin{lstlisting}
parts = read_metis_epart(path; one_based_parts = true)
\end{lstlisting}

\section{Visualize a partition}

Edit \code{NPROCS} in \code{examples/write\_partition\_vtk.jl} if required,
then run:

\begin{lstlisting}[language=bash]
julia --project=. examples/write_partition_vtk.jl
\end{lstlisting}

The script writes:

\begin{itemize}[leftmargin=*]
  \item one global legacy VTK file with \code{partition\_id} cell data;
  \item one owned-element VTK file per partition.
\end{itemize}

The principal visualization file is
\begin{center}
\api{examples/visualization/partitioned_mesh.vtk}.
\end{center}
Open it in ParaView, color by \code{partition\_id}, and select
``Surface With Edges''.

\section{Mesh-quality requirements}

\begin{itemize}[leftmargin=*]
  \item Tetrahedra must have four valid, distinct nodes.
  \item Every interior face must be shared by exactly two tetrahedra.
  \item Boundary triangles must match tetrahedral exterior faces.
  \item Volumes and face areas must be positive.
  \item The partition vector length must equal the tetrahedron count.
  \item Every partition ID must lie in \code{0:nranks-1}.
  \item Every MPI rank should normally own at least one element.
\end{itemize}

\chapter{Programming with DiscoGMPI}

\section{Construct a serial DG discretization}

\begin{lstlisting}
using DiscoGMPI

mesh = read_vtu_mesh("mesh.vtu")
dg = DGDiscretization(
    mesh,
    3;
    boundary_tag_name = "boundary_id",
    trace_tol = 1e-10,
    backend = SerialBackend(),
)
\end{lstlisting}

The resulting \code{DGDiscretization} contains the raw mesh, topology,
geometry, affine mappings, reference element, face operators, physical
operators, trace maps, flux faces, and backend.

\section{Select the threaded backend}

\begin{lstlisting}
dg = DGDiscretization(
    mesh,
    3;
    backend = ThreadedBackend(),
)
\end{lstlisting}

The threaded kernels use face coloring so concurrently processed faces do not
write to the same element residual.

\section{Create an initial Maxwell field}

\begin{lstlisting}
Efun = (x, y, z) -> (
    sin(pi * x),
    0.0,
    sin(pi * z),
)

Hfun = (x, y, z) -> (
    0.0,
    cos(pi * y),
    0.0,
)

U = interpolate_maxwell_field(mesh, dg.ref, Efun, Hfun)
\end{lstlisting}

Each component is stored as an \code{Np x Ne} matrix. Columns are elements and
rows are element-local DG nodes.

\section{Run a serial Hesthaven--Warburton case}

\begin{lstlisting}
registry = MaxwellBoundaryRegistry(
    Dict(10 => DiscoGMPI.MaxwellBC_PEC),
)
formulation = HesthavenWarburtonFormulation(MaxwellFlux_Upwind)

dt, _ = estimate_maxwell_dt(
    mesh,
    dg.geometry,
    dg.ref;
    CFL = 0.10,
)

run_maxwell_time_steps!(
    U,
    dg,
    registry,
    formulation;
    rk_order = 4,
    dt = dt,
    nsteps = 100,
    energy_every = 10,
)
\end{lstlisting}

\section{Run a Poisson-bracket case}

\begin{lstlisting}
formulation = PoissonBracketFormulation()

run_maxwell_partitioned_symplectic_time_steps!(
    U,
    dg,
    registry,
    formulation;
    psrk_order = 4,
    first_partition = :H,
    dt = dt,
    nsteps = 100,
    energy_every = 10,
)
\end{lstlisting}

\section{Build and run an MPI case}

\begin{lstlisting}
using MPI
using DiscoGMPI

MPI.Init()
comm = MPI.COMM_WORLD
rank = MPI.Comm_rank(comm)
nranks = MPI.Comm_size(comm)

root_mesh = rank == 0 ? read_vtu_mesh("mesh.vtu") : nothing
root_parts = rank == 0 ?
    read_metis_epart("mesh.mesh.epart.$nranks") :
    nothing

ddg = build_distributed_dg_from_root(
    root_mesh,
    root_parts,
    2;
    comm = comm,
)

Efun = (x, y, z) -> (sin(pi*x), 0.0, sin(pi*z))
Hfun = (x, y, z) -> (0.0, cos(pi*y), 0.0)
U = interpolate_maxwell_field(ddg, Efun, Hfun)

registry = MaxwellBoundaryRegistry(
    Dict(10 => DiscoGMPI.MaxwellBC_PEC),
)
formulation = HesthavenWarburtonFormulation(MaxwellFlux_Upwind)

run_distributed_maxwell_time_steps!(
    U,
    ddg,
    registry,
    formulation;
    rk_order = 3,
    dt = 1e-4,
    nsteps = 100,
    energy_every = 10,
)

MPI.Barrier(comm)
MPI.Finalize()
\end{lstlisting}

Launch it with:

\begin{lstlisting}[language=bash]
mpiexec -n 4 julia --project=. my_distributed_case.jl
\end{lstlisting}

\section{Distributed Poisson-bracket run}

\begin{lstlisting}
run_distributed_maxwell_partitioned_symplectic_time_steps!(
    U,
    ddg,
    registry,
    PoissonBracketFormulation();
    psrk_order = 4,
    first_partition = :H,
    dt = 1e-4,
    nsteps = 100,
    energy_every = 10,
)
\end{lstlisting}

\section{Energy diagnostics}

\begin{lstlisting}
serial_energy = maxwell_energy(
    U,
    dg.ref,
    dg.mappings,
)

global_energy = distributed_maxwell_energy(
    U,
    ddg,
)

println(global_energy.electric)
println(global_energy.magnetic)
println(global_energy.total)
\end{lstlisting}

Distributed energy is collective: every rank must call it in the same order.

\section{Global-to-local test initialization}

\code{localize\_maxwell\_field} extracts owned and halo columns from a global
field:

\begin{lstlisting}
local_U = localize_maxwell_field(global_U, ddg)
\end{lstlisting}

This function is primarily useful for verification against a serial reference.
For scalable production initialization, use
\code{interpolate\_maxwell\_field(ddg, Efun, Hfun)} so no global field is
required on worker ranks.

\chapter{Running Included Examples}

\section{Distributed Maxwell demonstration}

\begin{lstlisting}[language=bash]
mpiexec -n 2 julia --project=. examples/distributed_maxwell.jl
\end{lstlisting}

The script:
\begin{itemize}[leftmargin=*]
  \item builds a PEC cube only on rank zero;
  \item partitions its tetrahedra round-robin;
  \item distributes owned elements and face halos;
  \item initializes a cavity field locally;
  \item uses the Hesthaven--Warburton upwind formulation;
  \item advances two RK3 steps;
  \item reports ownership, neighbors, and global energy.
\end{itemize}

\section{Stationary PEC cube}

\begin{lstlisting}[language=bash]
julia --project=. examples/solver/stationary_pec_cube.jl \
  --cells=3 --order=2 --nsteps=20 --rk=4 --cfl=0.05 \
  --output=output/stationary_hw
\end{lstlisting}

The case uses the cube $[-1,1]^3$, PEC walls with boundary tag 10, and a
cavity mode with $\omega=\sqrt{3}\pi$. It prints energy, right-hand-side, and
tangential-electric-field diagnostics and writes final VTU data.

Display its complete options with:

\begin{lstlisting}[language=bash]
julia --project=. examples/solver/stationary_pec_cube.jl --help
\end{lstlisting}

\section{Stationary central-flux ESPRK case}

\begin{lstlisting}[language=bash]
julia --project=. examples/solver/stationary_pec_cube_esprk.jl \
  --cells=3 --order=2 --nsteps=20 --esprk=4 --first=H \
  --cfl=0.05 --output=output/stationary_esprk
\end{lstlisting}

This case combines the Hesthaven--Warburton central flux with an explicit
partitioned symplectic RK method. Both electric-first and magnetic-first stage
orders are accepted.

\section{Stationary Poisson-bracket ESPRK case}

\begin{lstlisting}[language=bash]
julia --project=. --threads=8 \
  examples/solver/stationary_pec_cube_pb_esprk.jl \
  --cells=3 --order=2 --nsteps=20 --esprk=4 \
  --backend=threaded --cfl=0.05 \
  --output=output/stationary_pb
\end{lstlisting}

The Poisson-bracket case requires \code{first=H}. It supports both the serial
and threaded backend options.

\section{Upwind convergence study}

\begin{lstlisting}[language=bash]
julia --project=. examples/solver/convergence_hw_upwind.jl \
  --cells=2,3,4,5 \
  --orders=1,2,4 \
  --time=0.005 \
  --cfl=0.1 \
  --rk=5 \
  --backend=serial \
  --output=output/convergence_hw_upwind.csv
\end{lstlisting}

For a threaded run:

\begin{lstlisting}[language=bash]
julia --project=. --threads=8 \
  examples/solver/convergence_hw_upwind.jl \
  --cells=2,3,4,5 --orders=1,2,4 --backend=threaded
\end{lstlisting}

The driver requires exactly four mesh levels. It creates deterministic
jittered tetrahedral cube meshes, computes continuous error norms, estimates
convergence rates, and writes CSV output.

\section{Central-flux partitioned convergence study}

\begin{lstlisting}[language=bash]
julia --project=. \
  examples/solver/convergence_hw_central_psrk_hfirst.jl \
  --cells=2,3,4,5 \
  --orders=1,2,4 \
  --psrk=2 \
  --first=H \
  --output=output/convergence_hw_central_psrk_hfirst.csv
\end{lstlisting}

\section{Poisson-bracket threaded convergence study}

\begin{lstlisting}[language=bash]
julia --project=. --threads=8 \
  examples/solver/convergence_pb_esprk_threaded.jl \
  --cells=3,4,5,6 \
  --orders=1,2,3,4 \
  --seconds=10 \
  --backend=threaded \
  --output=output/convergence_pb_esprk_threaded.csv
\end{lstlisting}

For this validation, the ESPRK order is fixed to spatial order plus one. A
second CSV file with \code{\_3dec} inserted before the extension is also
written.

\chapter{Testing}

\section{Package tests}

Run the non-MPI test target with:

\begin{lstlisting}[language=bash]
julia --project=. -e 'using Pkg; Pkg.test()'
\end{lstlisting}

The package suite covers loading, physical weak derivatives, Maxwell volume
and surface operators, formulations, backend dispatch, time integrators,
energy diagnostics, and distributed-mesh consistency constructed without an
MPI launcher.

\section{MPI face permutation test}

\begin{lstlisting}[language=bash]
mpiexec -n 2 julia --project=. test/test_mpi_face_permutation.jl
\end{lstlisting}

This test uses two tetrahedra whose shared face has different local vertex
ordering. It verifies all face-node values after MPI exchange, including
quadratic edge nodes.

\section{MPI METIS partition tests}

\begin{lstlisting}[language=bash]
mpiexec -n 2 julia --project=. test/test_mpi_metis_partition.jl
mpiexec -n 4 julia --project=. test/test_mpi_metis_partition.jl
\end{lstlisting}

These tests check owned and ghost counts, symmetric neighbor tables,
deterministic communication ordering, interface orientation, and generic trace
exchange.

\section{Distributed Maxwell equivalence tests}

\begin{lstlisting}[language=bash]
mpiexec -n 2 julia --project=. test/test_mpi_distributed_maxwell.jl
mpiexec -n 4 julia --project=. test/test_mpi_distributed_maxwell.jl
\end{lstlisting}

The two-rank test compares distributed and serial results for:

\begin{itemize}[leftmargin=*]
  \item Hesthaven--Warburton central flux;
  \item Hesthaven--Warburton upwind flux;
  \item Poisson-bracket formulation;
  \item exchanged ghost traces;
  \item zero ghost residuals;
  \item global energy;
  \item one explicit RK3 step;
  \item one partitioned symplectic RK2 step.
\end{itemize}

The four-rank chain test exercises ranks with two neighbors and verifies
upwind right-hand sides, traces, ghost residuals, and RK3 stepping. The file
also verifies that root input errors propagate collectively without an MPI
deadlock.

\section{Recommended pre-commit sequence}

\begin{lstlisting}[language=bash]
julia --project=. -e 'using Pkg; Pkg.test()'
mpiexec -n 2 julia --project=. test/test_mpi_face_permutation.jl
mpiexec -n 2 julia --project=. test/test_mpi_metis_partition.jl
mpiexec -n 4 julia --project=. test/test_mpi_metis_partition.jl
mpiexec -n 2 julia --project=. test/test_mpi_distributed_maxwell.jl
mpiexec -n 4 julia --project=. test/test_mpi_distributed_maxwell.jl
mpiexec -n 2 julia --project=. examples/distributed_maxwell.jl
\end{lstlisting}

\chapter{Benchmarking}

\section{Single benchmark}

\begin{lstlisting}[language=bash]
julia --project=. --threads=8 \
  benchmark/solver/benchmark_backends.jl \
  --formulation=hw-upwind \
  --cells=6 \
  --order=3 \
  --seconds=3
\end{lstlisting}

Select the Poisson-bracket operator with
\code{--formulation=pb}. The benchmark reports minimum, median, and mean
times, allocations, memory, speedup, efficiency, and serial-versus-threaded
right-hand-side error.

\section{Thread scaling sweep}

\begin{lstlisting}[language=bash]
THREAD_COUNTS="1 2 4 8" \
  bash benchmark/solver/run_thread_scaling.sh \
  --formulation=hw-upwind --cells=6 --order=3
\end{lstlisting}

Outputs:

\begin{itemize}[leftmargin=*]
  \item \code{benchmark/solver/thread\_scaling\_raw.txt};
  \item \code{benchmark/solver/thread\_scaling.csv}.
\end{itemize}

\section{Plot scaling data}

\begin{lstlisting}[language=bash]
python3 benchmark/solver/plot_thread_scaling.py
\end{lstlisting}

The script writes:

\begin{itemize}[leftmargin=*]
  \item \code{benchmark/solver/thread\_speedup.png};
  \item \code{benchmark/solver/thread\_efficiency.png}.
\end{itemize}

Install Matplotlib separately if it is not available:

\begin{lstlisting}[language=bash]
python3 -m pip install matplotlib
\end{lstlisting}

\section{MPI performance measurements}

The current repository does not contain a full MPI scaling harness. For
meaningful measurements:

\begin{enumerate}[leftmargin=*]
  \item exclude Julia compilation and mesh distribution from the timed region;
  \item warm up at least one right-hand-side evaluation;
  \item place an \code{MPI.Barrier} immediately before timing;
  \item reduce the maximum elapsed time across ranks;
  \item report owned elements, interface faces, polynomial order, and
        $N_{fp}$;
  \item separate right-hand-side, trace exchange, energy reduction, and output
        costs;
  \item compare strong and weak scaling.
\end{enumerate}

\chapter{Output and Visualization}

\section{Serial VTU field output}

The stationary PEC examples contain helper functions that write either:

\begin{itemize}[leftmargin=*]
  \item cell-centered fields for order zero; or
  \item discontinuous vertex fields for order one and above.
\end{itemize}

The discontinuous writer duplicates tetrahedron vertices per element. This is
intentional: a DG solution may have two different traces at the same geometric
vertex.

\section{Partition output}

\api{write_partition_vtk} writes the complete mesh with a partition cell
array. \api{write_partition_piece_vtks} writes one owned-element mesh per
partition. These are mesh visualization utilities; they do not write the
distributed Maxwell solution.

\section{Distributed solution output}

A general MPI solution gather or parallel VTK writer is not yet part of the
public API. A case-specific writer should:

\begin{enumerate}[leftmargin=*]
  \item write owned elements only;
  \item preserve global element IDs;
  \item avoid writing ghost elements as physical duplicates;
  \item generate one file per rank or use a parallel VTK metadata file;
  \item represent DG discontinuities explicitly, normally by duplicating
        element-local points.
\end{enumerate}

\chapter{Diagnostics and Troubleshooting}

\section{MPI run hangs}

\begin{itemize}[leftmargin=*]
  \item Confirm every rank enters the same collective constructor and time
        integration sequence.
  \item Confirm the launch rank count matches the partition file.
  \item Confirm MPI.jl and \code{mpiexec} use compatible MPI libraries.
  \item Do not let only rank zero call \code{distributed\_maxwell\_energy};
        it contains an \code{Allreduce}.
  \item Ensure all ranks use the same communication tag for a given exchange.
\end{itemize}

\section{Wrong interface flux or rank-dependent error}

\begin{itemize}[leftmargin=*]
  \item Run \code{test\_mpi\_face\_permutation.jl}.
  \item Check tetrahedron connectivity and face orientation.
  \item Check that face-node generation is identical on all ranks.
  \item Tighten or relax \code{trace\_tol} only after inspecting mesh scale.
  \item Verify that the partition includes complete face adjacency, not merely
        node adjacency.
\end{itemize}

\section{Missing or incorrect boundary conditions}

\begin{itemize}[leftmargin=*]
  \item Confirm the VTU file contains triangular boundary cells.
  \item Confirm their cell-data array is named \code{boundary\_id}, or pass a
        different \code{boundary\_tag\_name}.
  \item Print \code{unique(tri\_data(mesh, "boundary\_id"))}.
  \item Register every physical PEC tag explicitly.
  \item Remember that tag 10 is a convention in the examples, not a hard
        requirement of the solver.
\end{itemize}

\section{Unstable time integration}

\begin{itemize}[leftmargin=*]
  \item Reduce the CFL factor.
  \item Inspect the worst element reported by
        \code{element\_size\_diagnostics}.
  \item Check for degenerate or highly skewed tetrahedra.
  \item Verify positive $\varepsilon$ and $\mu$.
  \item Use the upwind Hesthaven--Warburton flux when damping unresolved jumps
        is acceptable.
  \item Confirm that the chosen RK or ESPRK order is supported.
\end{itemize}

\section{Poisson-bracket errors}

\begin{itemize}[leftmargin=*]
  \item Use \code{PoissonBracketFormulation()} with central flux only.
  \item Use the partitioned symplectic driver.
  \item Set \code{first\_partition = :H}.
  \item Use ESPRK order one through six.
\end{itemize}

\section{Threaded run is slower}

\begin{itemize}[leftmargin=*]
  \item Increase the problem size; small meshes are overhead dominated.
  \item Verify the actual Julia thread count.
  \item Avoid oversubscribing cores in hybrid MPI/threaded runs.
  \item Pin ranks and threads appropriately for the machine.
  \item Use the benchmark harness and compare median times.
\end{itemize}

\section{Memory use is larger than expected}

\begin{itemize}[leftmargin=*]
  \item Use \code{build\_distributed\_dg\_from\_root}, not the replicated
        constructor.
  \item Remember that each rank stores owned elements plus one face halo.
  \item Higher polynomial order increases storage in proportion to $N_p$.
  \item Explicit RK methods store several complete right-hand sides.
  \item Root temporarily retains the complete global mesh and partition during
        distribution.
\end{itemize}

\chapter{Current Limitations and Roadmap}

\section{Current limitations}

\begin{itemize}[leftmargin=*]
  \item Three-dimensional tetrahedral meshes only.
  \item Affine straight-sided tetrahedral geometry.
  \item Source-free Maxwell equations.
  \item Scalar, run-wide $\varepsilon$ and $\mu$ in the Maxwell operator.
  \item PEC is the implemented physical boundary condition for the MPI
        Maxwell path.
  \item Distributed periodic boundaries are not implemented.
  \item No PML.
  \item No material-interface Riemann solver for discontinuous impedance.
  \item No adaptive mesh refinement.
  \item No distributed mesh repartitioning during a run.
  \item No general parallel field-output layer.
  \item Root-only loading removes global mesh replication from workers but
        root still reads and partitions the complete mesh.
  \item Communication is nonblocking, but currently completed before local DG
        evaluation; volume work is not yet overlapped with trace exchange.
  \item Stored ghost elements currently pass through local volume evaluation
        before ghost residuals are cleared.
\end{itemize}

\section{Recommended next development stages}

\begin{enumerate}[leftmargin=*]
  \item Add distributed periodic face pairing.
  \item Add sources and spatially varying materials.
  \item Add material-jump fluxes using side-dependent impedance and
        admittance.
  \item Add PML.
  \item Split the right-hand side into owned volume work and interface surface
        work, then overlap communication.
  \item Skip ghost volume computation.
  \item Add parallel solution output.
  \item Add MPI strong- and weak-scaling benchmarks.
  \item Replace root assembly with direct per-rank partition-file loading for
        meshes too large for one rank.
\end{enumerate}

\appendix

\chapter{API Quick Reference}

\begin{longtable}{p{0.45\textwidth}p{0.47\textwidth}}
\toprule
API & Purpose \\
\midrule
\endhead
\api{read_vtu_mesh(path)} & Read XML VTU tetrahedra, boundary triangles, and cell
data into \code{RawVTUMesh}. \\
\api{DGDiscretization(mesh,order;...)} & Build complete serial/threaded DG
operators. \\
\api{SerialBackend()} & Select scalar local kernels. \\
\api{ThreadedBackend()} & Select Julia-threaded local Maxwell kernels. \\
\api{HesthavenWarburtonFormulation(kind)} & Strong-form Maxwell DG with central or
upwind flux. \\
\api{PoissonBracketFormulation()} & Partitioned weak Maxwell operator with central
flux. \\
\api{MaxwellBoundaryRegistry(dict)} & Map integer boundary tags to boundary kinds. \\
\api{interpolate_maxwell_field(...)} & Sample analytic electric and magnetic
functions at DG nodes. \\
\api{estimate_maxwell_dt(...)} & Compute the volume-to-surface CFL estimate. \\
\api{maxwell_energy(...)} & Compute serial electromagnetic energy. \\
\api{run_maxwell_time_steps!(...)} & Advance Hesthaven--Warburton with explicit
RK. \\
\api{run_maxwell_partitioned_symplectic_time_steps!(...)} & Advance a
partitioned Maxwell system with ESPRK. \\
\api{read_metis_epart(path)} & Read zero-based element partition IDs. \\
\api{build_distributed_mesh(...)} & Build one rank's owned-plus-halo mesh from
global arrays. \\
\api{build_distributed_dg(...)} & Build distributed DG when the global mesh is
replicated. \\
\api{build_distributed_dg_from_root(...)} & Distribute rank-local DG meshes from
a root-owned global mesh or root-read files. \\
\api{exchange_maxwell_ghost_traces!(...)} & Exchange six Maxwell face traces. \\
\api{distributed_maxwell_energy(...)} & Compute owned-only local energy followed
by a global sum. \\
\api{run_distributed_maxwell_time_steps!(...)} & Distributed explicit RK
driver. \\
\api{run_distributed_maxwell_partitioned_symplectic_time_steps!(...)} &
Distributed Poisson-bracket ESPRK driver. \\
\api{write_metis_mesh_from_vtk(...)} & Convert a supported legacy VTK tetra mesh
to METIS format. \\
\api{write_partition_vtk(...)} & Write a legacy VTK mesh colored by partition. \\
\api{write_partition_piece_vtks(...)} & Write owned-element mesh pieces. \\
\bottomrule
\end{longtable}

\chapter{Command Quick Reference}

\begin{lstlisting}[language=bash]
# Instantiate
julia --project=. -e 'using Pkg; Pkg.instantiate()'

# Package tests
julia --project=. -e 'using Pkg; Pkg.test()'

# Distributed Maxwell demo
mpiexec -n 2 julia --project=. examples/distributed_maxwell.jl

# MPI tests
mpiexec -n 2 julia --project=. test/test_mpi_face_permutation.jl
mpiexec -n 2 julia --project=. test/test_mpi_metis_partition.jl
mpiexec -n 4 julia --project=. test/test_mpi_metis_partition.jl
mpiexec -n 2 julia --project=. test/test_mpi_distributed_maxwell.jl
mpiexec -n 4 julia --project=. test/test_mpi_distributed_maxwell.jl

# Mesh inspection
julia --project=. examples/solver/read_vtu_mesh.jl mesh.vtu

# METIS preparation and partitioning
julia --project=. examples/write_metis.jl
mpmetis examples/meshes/tet_mesh.mesh 4

# Partition visualization
julia --project=. examples/write_partition_vtk.jl

# Serial stationary case
julia --project=. examples/solver/stationary_pec_cube.jl --help

# Threaded Poisson-bracket case
julia --project=. --threads=8 \
  examples/solver/stationary_pec_cube_pb_esprk.jl --help

# Backend benchmark
julia --project=. --threads=8 \
  benchmark/solver/benchmark_backends.jl --help

# Thread scaling
THREAD_COUNTS="1 2 4 8" \
  bash benchmark/solver/run_thread_scaling.sh

# Plot benchmark data
python3 benchmark/solver/plot_thread_scaling.py

# Build this guide
pdflatex docs/DiscoGMPI_User_Guide.tex
pdflatex docs/DiscoGMPI_User_Guide.tex
\end{lstlisting}

\chapter{Reproducibility Checklist}

Record the following information with every reported result:

\begin{itemize}[leftmargin=*]
  \item \pkg{} commit hash and working-tree status;
  \item Julia version and project/manifest state;
  \item MPI implementation and version;
  \item number of MPI ranks and Julia threads per rank;
  \item CPU/node allocation and affinity settings;
  \item mesh file, tetrahedron count, boundary tags, and partition file;
  \item polynomial order $N$, $N_p$, and $N_{fp}$;
  \item formulation and flux;
  \item RK or ESPRK order and partition order;
  \item $\varepsilon$, $\mu$, CFL, $\Delta t$, and number of steps;
  \item initial condition and boundary registry;
  \item initial/final energy and error norms;
  \item output and benchmark methodology.
\end{itemize}

\backmatter

\chapter*{License}
\addcontentsline{toc}{chapter}{License}

\pkg{} is distributed under the MIT License. See \code{LICENSE.md} in the
repository for the complete license text.

\end{document}
